\section{Introduction}
\label{sec:intro}

A language agent that persists across sessions must decide what to do with each
candidate memory: retain the current record, replace it with a newer one, or
decline to commit because the evidence is insufficient. Deployed systems often
make this decision through opaque LLM tool calls
\citep{packer2023memgpt,chhikara2025mem0} or additive event logs
\citep{park2023generative}. A wrong write can remain latent until much later: an
ambiguous statement becomes a durable fact, a stale preference survives, or a
correction erases its predecessor. End-to-end QA then reveals an error without
identifying which layer caused it.

We take the write transition itself as the object of evaluation. The action is
\keep, lineage-preserving \supersede, or write-time \abstain{} over typed,
provenance-tracked records. This combines established ideas from selective
prediction \citep{chow1970reject,geifman2017selective}, belief revision
\citep{alchourron1985agm}, and agent memory operations
\citep{yan2025memoryr1,chhikara2025mem0,wang2026toki}. We claim neither the
individual operations nor abstention, typing, provenance, and lineage as new.
The contribution is a harness that makes their decision-, state-, and
answer-layer consequences separately observable.

We first exercise the harness on \mgp, a $180$-item, English, synthetic pilot.
The collection was authored by the same project that implemented the reference
policy and has only one independent human re-annotation of its test split. It is
therefore a diagnostic instrument, not an independently established benchmark.
Rather than treating a held-out item split as sufficient, we explicitly test
whether scenario metadata predicts the action and whether item-level inference
understates within-template dependence.

Our contributions and findings are:
\begin{enumerate}[leftmargin=1.5em,topsep=2pt,itemsep=1pt]
\item We formalize the three write actions and release a layered evaluator for
action correctness, state invariants, lineage, stale retrieval, selective risk,
and asymmetric write costs.
\item We add train/dev-fitted shortcut baselines and template-family/minimal-pair
cluster bootstrap. Conflict type alone predicts every held-out test action;
type alone predicts $96.7\%$. This is direct evidence that the current pilot
does not identify general evidence reasoning.
\item On the held-out test split, a transparent reference policy reaches
$0.8889$ accuracy and a structured LLM policy reaches $0.9000$; the paired
family-cluster interval includes zero. Zero observed reference-policy false
commits has a non-zero finite-sample upper bound and trades against missed
updates, so the preferred policy changes with the cost matrix.
\end{enumerate}

This reframing turns a weakness into a falsifiable result: an item-level test
split can look clean while a family label still determines the answer. The
released artifacts make that failure visible and specify the gate for a
successor collection: independent construction, all three actions within each
conflict family, blinded multi-human annotation, trajectory-level cases, and a
hidden test set frozen after policy development.
